 \setlength{\baselineskip}{20pt}
\begin{englishabstract}
	With the rapid development of automated parking and intelligent driving systems, parking slot detection has become a critical component of environment perception, where detection accuracy and robustness directly affect system safety and reliability. Unlike conventional object detection tasks, parking slot detection is characterized by strong structural constraints, high annotation costs, and long-tailed scene distributions, which pose significant challenges to existing methods in complex real-world environments. To address these issues, this dissertation conducts a systematic study on parking slot detection from three perspectives: large-scale data construction, structure-aware end-to-end modeling, and multi-frame temporal enhancement.

First, a large-scale real-world parking slot detection dataset, termed CRPS-D, is constructed, consisting of 35,379 images with diverse scenarios, high parking slot density, and a large proportion of abnormal slots. Based on this dataset and the sparse, structure-constrained nature of parking slot keypoints, a teacher–student-based semi-supervised learning framework is proposed. By incorporating confidence-guided mask consistency (CGM) and an adaptive feature perturbation mechanism (Adaptive-VAT), the framework effectively exploits spatial structural priors from unlabeled data. Experimental results demonstrate that, using only one-twelfth of the labeled training data, the proposed method achieves an average precision (AP) of 90.47\% on the ps2.0 dataset and 79.75\% on CRPS-D, significantly outperforming fully supervised baselines trained under the same annotation budget.

Second, a structure-aware end-to-end detection framework for single-frame parking slot detection is presented. To overcome the limitations of heuristic post-processing and rule-based matching in existing approaches, geometric topological constraints are explicitly embedded into the network to jointly optimize keypoint localization and parking slot structure inference. This end-to-end formulation substantially improves structural consistency, particularly in dense and partially occluded scenarios. Experimental results show that the proposed method achieves 99.37\% precision, 99.47\% recall, and 99.47\% AP on the ps2.0 dataset, and attains 92.21\% precision, 85.54\% recall, and 83.58\% AP on the more challenging CRPS-D dataset, surpassing state-of-the-art end-to-end methods.

Finally, a multi-frame temporal feature fusion approach is proposed to further enhance detection robustness. By extending static spatial modeling to spatio-temporal modeling, the method exploits temporal correlations in consecutive video frames to recursively refine and smooth structure-aware features, improving temporal consistency under occlusion and environmental noise. Experiments on the self-collected PSTD temporal dataset demonstrate that the proposed approach achieves 92.53\% precision, 92.33\% recall, and 91.97\% AP. Moreover, the model maintains a stable inference speed of approximately 30 FPS on embedded automotive platforms, satisfying the real-time requirements of practical automated parking systems.
	
	\noindent\englishkeywords{Parking Slot Detection; Dataset; Semi-supervised Learning; Keypoint Detection; Structure-aware Modeling; Temporal Modeling}
\end{englishabstract}